\section{Features, profiles and the strategic prior}
\label{sec:features}

\subsection{Fixed-dimensional vectors for variable objects}
The state vector has $d_s=384$ coordinates and each candidate's action vector
has $d_a=192$. A named feature $f$ with value $v_f$ is added by signed hashing:
\begin{equation}
  \phi_j \mathrel{+}=\sigma(f)\,v_f,\qquad
  j=H(f)\bmod(d-16),\qquad \sigma(f)\in\{-1,+1\}.
\end{equation}
Both $H$ and $\sigma$ come from an eight-byte BLAKE2b digest: the first four
bytes, little-endian, give $H$, and the parity of the fifth byte gives the
sign. The last sixteen coordinates are reserved for dense numbers. That
leaves 368 hashed state buckets and 176 hashed action buckets. Feature names
are namespaced by role and position: species, item, ability, HP, status,
boosts, stats, types, moves with their PP ratio, volatiles, field and side
conditions. The last sixteen public move events are added with weight
$1/\sqrt{i+1}$, counting back from the newest. A research prior adds
\code{meta/<species>} frequencies. It is built from the top 32 team-sheet
frequencies, overlaid by the top 32 usage entries. Dense state slots hold turn
progress $\min(\text{turn}/50,2)$, total own HP divided by six, and
$|\A|/1000$.

The whole vector is clipped to $[-5,5]$. After clipping, the decision code
adds the scout's top-three move probabilities
(Section~\ref{sec:scout}). The \emph{stored} state therefore can leave the
clipped interval.

\begin{proposition}[Hashing loses distinctions]
\label{prop:hash}
A signed hash map from $m$ categorical indicator coordinates into $d'<m$
buckets ($d'=d-16$ here) is not injective on $\R^m$. Any later function of the
hashed vector, including clipping, cannot restore the lost distinctions.
\end{proposition}
\begin{proof}
The map is linear with rank at most $d'<m$, so its null space is nontrivial.
Two inputs that differ by a null-space vector hash to the same output, and
any later function maps equal inputs to equal outputs.
\end{proof}
This is a capacity statement, not a claim that two particular battle states
collide. Expected-collision results for feature hashing assume random hash
functions \citep{featurehashing}; a fixed digest over a realized vocabulary
does not inherit them automatically.

\subsection{Feature profiles}
A checkpoint stores a \code{feature_profile}. Training excludes episodes
recorded under any other profile, so a profile is a versioned input semantics.
The schema number alone is not enough to identify the inputs. Nineteen profiles
exist; Table~\ref{tab:profiles} lists the families.

\begin{table}[htbp]
\centering\small
\begin{tabularx}{\linewidth}{>{\raggedright\arraybackslash}p{0.27\linewidth} X}
\toprule
Profile family & What it changes \\
\midrule
\code{legacy}, \code{weather-v1} & Original features and the power/type prior
of Section~\ref{sec:legacyprior}; weather-aware move data. \\
\code{tactics-v1}, \code{rain-v1}, \code{mechanics-v1} & Damage-based turn
priors and rain-team mechanics. \code{tactics-v1} and \code{rain-v1}, with
\code{opening-v1}, are the only profiles that use the corrected coloured-HP
parser. \\
\code{preview-*}, \code{lineup-v1}, \code{opening-*}, \code{mega-v1} &
Preview priors (coverage, lead pressure, exposure) and exact prospective Mega
stats under Champions rules. \\
\code{strategic-v1} (\textbf{deployed}) & Scenario-evaluation prior for turns
and previews (Section~\ref{sec:strategic}); bench canonicalization. \\
\code{strategic-v2..v5} & More mechanics, preview plan summaries, Mega
uncertainty. v2--v4 were evaluated and rejected; the v5 experiment was paused
before a completed gate (Section~\ref{sec:evidence}). \\
\bottomrule
\end{tabularx}
\caption{Feature-profile families. Only \code{strategic-v1} is deployed.}
\label{tab:profiles}
\end{table}

\subsection{The deployed strategic prior}
\label{sec:strategic}
The final action coordinate holds a scalar prior $q_0(x,a)$. Under
\code{strategic-v1} it comes from a one-turn \emph{scenario evaluation} in
\code{ml/strategy.py}, which has four steps.

\paragraph{Opponent hypotheses.}
For each live opposing active Pok\'{e}mon, candidate moves are drawn from the
disclosed sheet if there is one. Otherwise they are the observed moves plus
the five most frequent executed moves for the species in the checkpoint's
\code{strategy_knowledge}. That table holds counts from 766 earlier games.
Each move gets the weight
\begin{equation}
  w=2+\sqrt{\text{count}}+\min(3,\text{recent}),
\end{equation}
with Protect multiplied by 1.5 below 40\% HP, by 0.65 otherwise, and by 0.3
if it was used on the previous turn. The two strongest attacks and one support
move are kept. A single-target attack is split over targets: weight~1 goes to
the target with the largest damage relative to its remaining HP
($\text{damage}/\max(0.2,h)$), and 0.45 to the others. Two further hypotheses are
added when they apply:
\begin{itemize}
  \item a switch-in from the revealed bench, with weight $0.12\sum w$;
  \item a 50/50 Mega-forme split, when the stone is disclosed and the side's
    Mega is unspent.
\end{itemize}

\paragraph{Joint scenarios.}
Per-position hypotheses are multiplied across the two opposing positions.
Impossible combinations, such as two Megas or two switches into one reserve,
are removed. The six most probable scenarios are kept and renormalized to
probabilities $p_k$.

\paragraph{Deterministic resolution.}
For our candidate $a$ and scenario $k$, \code{simulate} resolves one turn in
priority and speed order, including Trick Room. It accounts for Protect, Wide
Guard, redirection, Helping Hand, Fake Out, weather, terrain, Intimidate on
entry, drain, recoil and Focus Sash. Damage uses the standard formula
\begin{equation}
  \text{dmg}=\Big(\tfrac{22\,\mathrm{BP}\,A/D}{50}+2\Big)\cdot\mathrm{STAB}\cdot E\cdot
  \text{burn}\cdot\text{weather}\cdot0.75^{[\text{spread}]}\cdot0.925\cdot\text{acc}\,/\,\text{maxHP},
\end{equation}
where unknown stats are approximated as base~$+\,35$ with stage multipliers, and
unknown maximum HP as base HP~$+\,90$. The constant 0.925 is the mean damage
roll. The outcome value is
\begin{equation}
  v_k(a)=\sum_{\text{foes}}\big[\Delta h+0.8\,\mathbf 1_{\mathrm{KO}}\big]
  -\sum_{\text{own}}\big[\Delta h(1+0.2\rho)+(0.9+\rho)\mathbf 1_{\mathrm{KO}}\big]
  -0.08\,\#\text{Protect}-0.10\,\#\text{switch}+\dots,
\end{equation}
where $\rho\in[0,0.5]$ is a scarcity bonus for a Pok\'{e}mon that is the only
good answer to some opponent. Smaller terms reward a switch-in's next-turn
threat and penalize spending the Mega.

\paragraph{Risk-sensitive aggregation and squashing.}
If speeds tie, each $v_k$ averages both orders. The turn score is
\begin{equation}
  s(a)=3\Big[(1-\kappa)\sum_kp_kv_k(a)+\kappa\min_kv_k(a)\Big],\qquad
  \kappa=\clip\big(0.15+0.05(R_{\mathrm{own}}-R_{\mathrm{foe}}),0.05,0.25\big),
  \label{eq:turnscore}
\end{equation}
where $R_{\mathrm{own}}$ sums our HP fractions. $R_{\mathrm{foe}}$ sums the
revealed foes' HP fractions and counts each unseen foe as one. Preview
candidates use \code{opening_score} instead. It takes the better of no-Mega
and one-Mega modes over coverage, lead pressure and exposure, adds terms for
rain dependency, Intimidate, Fake Out and Trick Room, and doubles the total.
Finally,
\begin{equation}
  q_0(a)=g(s(a)),\qquad g(s)=\frac{4s}{4+|s|}.
  \label{eq:squash}
\end{equation}

\begin{proposition}[Properties of the risk aggregate and the squash]
\label{prop:squash}
\leavevmode
\begin{enumerate}
  \item For any scenario values,
    $\min_kv_k\le(1-\kappa)\sum_kp_kv_k+\kappa\min_kv_k\le\sum_kp_kv_k$.
    The aggregate is monotone in each $v_k$, and adding a constant $c$ to every
    $v_k$ adds $c$ to it.
  \item The function $g$ is odd, strictly increasing and a bijection from $\R$
    onto $(-4,4)$. It satisfies $0<g'(s)=16/(4+|s|)^2\le1$, so it is
    1-Lipschitz. Its ranking of candidates equals the ranking by $s$.
\end{enumerate}
\end{proposition}
\begin{proof}
(1) Since $\min_kv_k\le\sum_kp_kv_k$, a convex combination of the two lies
between them. Each part is monotone and translation-equivariant, and so is
their combination. (2) For $s\ge0$, $g(s)=4s/(4+s)$ has derivative
$16/(4+s)^2>0$; oddness extends this to $s<0$. The function is continuous at
0, and $g(s)\to\pm4$ as $s\to\pm\infty$. The derivative is largest at $s=0$,
where it equals~1, which gives the Lipschitz bound. A strictly increasing map
preserves order.
\end{proof}

\intuition{The squash replaces a hard clip, which had flattened 288 of the 360
preview candidates into one tied value (\code{docs/strategic-validation.md}).
The pessimism weight $\kappa$ \emph{increases} when we are ahead on resources:
the prior becomes more risk-averse when there is more to lose. The final
\code{clip} to $[-4,4]$ in the source is redundant after $g$.}

\begin{example}[A strategic prior, end to end]
\label{ex:strategic}
Three scenarios have values $(0.9,0.4,-0.6)$ and probabilities
$(0.5,0.3,0.2)$, with equal resources, so $\kappa=0.15$. The expectation is
$0.45$, the minimum is $-0.6$, and
$s=3[0.85(0.45)+0.15(-0.6)]=0.8775$. The squash gives
$q_0=4(0.8775)/4.8775\approx0.720$. At the deployed temperature $\tau=0.25$
(Section~\ref{sec:network}), this prior adds $0.720/0.25\approx2.88$ to the
effective logit before any learned residual.
\end{example}

\subsection{The legacy power prior}
\label{sec:legacyprior}
Earlier profiles used the rough pressure term
\begin{equation}
  p(m,j)=\frac{\mathrm{BP}(m)}{100}\,\mathrm{acc}(m)\,\mathrm{STAB}(m)\,E(m,j),
\label{eq:prior}
\end{equation}
summed over live targets, negated against allies, minus 0.4 per switch, plus
0.15 for Protect, Detect or Wide Guard, and clipped to $[-4,4]$.

\begin{example}[A legacy prior]
Base power 80, accuracy 1, STAB 1.5 and a super-effective multiplier of 2 give
$0.8\cdot1.5\cdot2=2.4$. Paired with Protect the candidate scores 2.55; paired
with a switch it scores 2.0. These numbers are preferences, not HP damage or
win probabilities. At preview this legacy prior is zero; the strategic prior
is not.
\end{example}

The legacy term still feeds a dense action slot, but it is not the deployed
prior. Neither prior models full speed control, stat spreads, most items, or
the opponent's best response. Search (Section~\ref{sec:search}) addresses
exactly that last omission.

\source{\code{Features.add/state/action} in \code{ml/features.py};
\code{hypotheses}, \code{prepare}, \code{simulate}, \code{turn_score} and
\code{opening_score} in \code{ml/strategy.py}; \code{damage} in
\code{ml/tactics.py}; \code{preview_score} in \code{ml/opening.py};
\code{species_prior} in \code{ml/research.py}.}
